% GENERATED FILE. Edit canonical lessons, then run npm run generate:books.
% canonical-source-hash: fnv1a64:d507774498de4333
% canonical-lessons: HI-C01-namaste, HI-C01-namaskar, HI-C01-dhanyavad, HI-C01-shukriya, HI-C01-alvida, HI-W01-shirorekha-na-ma, HI-W01-na-ma, HI-W02-abugida-ka-ta, HI-W02-ka-ta-mouth-order, HI-C01-practice

\chapter{Greetings}
\label{ch:greetings}
\hlchaptermodality{1}

\begin{chapteropening}
\textbf{By the end of this chapter:} \emph{I can greet, thank, and take my leave in Hindi, and write by hand the first Devanagari letters those words are built from.}

Greet, thank formally and casually, and take a weighty leave in Hindi from memory; name each word's heritage; then write \hi{न}, \hi{म}, \hi{क} and \hi{त} by hand, body first and bar last, and say what orders the stop families along the mouth.
\end{chapteropening}

\section[\texorpdfstring{\hi{नमस्ते} (namaste)}{नमस्ते (namaste)}]{Namaste — \textquotedblleft{}I bow to you\textquotedblright{}}

\label{lesson:HI-C01-namaste}



\begin{quote}
The first Hindi word, and the most important — a greeting that is literally a small act of reverence. You'll meet its first and last written pieces today; the middle comes next lesson.
\end{quote}

\subsection*{The letters in this word}
\emph{(If you already read Devanagari, skim this — it's here so the book needs no prior knowledge.)} Hindi is written \textbf{left to right}, and every consonant carries a built-in \textbf{\textquotedblleft{}a\textquotedblright{}} unless something changes it.

\begin{itemize}
\raggedright
  \item \textbf{\hi{न}} = \textquotedblleft{}na\textquotedblright{} — the word's first letter, with its built-in \emph{a}.
  \item \textbf{\hi{त}} = \textquotedblleft{}ta\textquotedblright{} — another consonant, with the same built-in \emph{a}.
  \item \textbf{\hi{े}} is a \textbf{vowel sign} (a \emph{mātrā}) that changes the built-in \emph{a} to \textquotedblleft{}e\textquotedblright{}: \textbf{\hi{त} + \hi{े} $\to$ \hi{ते}} = \textquotedblleft{}te.\textquotedblright{}
\end{itemize}

So the word begins with \textbf{\hi{न}} \emph{na} and ends with \textbf{\hi{ते}} \emph{te}. The middle — a \emph{ma} and a bare \emph{s} — comes in the next lesson; the whole written word follows later in this chapter.

Two facts already: consonants carry \emph{a}, and a \emph{mātrā} changes the vowel. The next lesson adds the third.

\begin{cousinweb}
\emph{Namaste} = \emph{namaḥ} (\textquotedblleft{}a bow, reverence, homage\textquotedblright{}) + \textbf{\hi{ते}} (\emph{te}, \textquotedblleft{}to you\textquotedblright{}) = literally \textbf{\textquotedblleft{}reverence to you,\textquotedblright{} \textquotedblleft{}I bow to you.\textquotedblright{}} The root is Sanskrit \emph{nam}, \textquotedblleft{}to bow, bend, pay homage\textquotedblright{} — the same root behind \emph{namaskār} (next lesson) and borrowed whole into English as \textbf{namaste} and \textbf{namaskar}.
\end{cousinweb}

\begin{culture}
\emph{namaste} is a genuine gesture, not a neutral \textquotedblleft{}hi\textquotedblright{}: it's said with the palms pressed together (the \emph{añjali} gesture), acknowledging the other person with respect. It works as both \textbf{hello and goodbye}, at any time of day, and it's polite/respectful without being stiff — safe with almost anyone. (Hindi also has warmer and more casual options you'll meet, and a whole second, Persian- flavored set of words — coming in a few lessons.)
\end{culture}

\begin{sounds}
Hindi reads \textbf{left to right}, and every consonant carries a built-in “a”: \hi{न} na, \hi{त} ta. A vowel sign changes the vowel: \hi{े} makes “e” (\hi{ते} = te). Say the four beats: \emph{na · ma · s · te}. The \emph{s} has no vowel of its own and leans straight into \emph{te}.
\end{sounds}

\subsection*{Your turn}
Say these aloud:

\begin{itemize}
\raggedright
  \item its beats — na · ma · s · te $\to$ \textquotedblleft{}namaste\textquotedblright{}
  \item read \hi{ते} — \textquotedblleft{}te\textquotedblright{} — and what the sign \hi{े} did to \hi{त} (changed \emph{a} to \emph{e})
  \item \textquotedblleft{}namaste,\textquotedblright{} and mean it — \textquotedblleft{}I bow to you\textquotedblright{}
\end{itemize}

\begin{tcolorbox}[breakable,colback=teal!4,colframe=teal!35!black,title={Before you move on}]
\emph{Read it:} \textbf{\hi{न}} and \textbf{\hi{ते}} (\emph{na}, \emph{te}.) What does \emph{namaste} literally mean? (\textquotedblleft{}I bow to you\textquotedblright{} — \emph{namaḥ} + \emph{te}.) What does a \emph{mātrā} do to a consonant? (Changes its built-in \emph{a}.)
\end{tcolorbox}
\section[\texorpdfstring{\hi{नमस्कार} (namaskār)}{नमस्कार (namaskār)}]{Namaskār — the more formal bow}

\label{lesson:HI-C01-namaskar}



\begin{quote}
The dressier sibling of \emph{namaste} — same root, one new piece, and the letters that finish \emph{namaste} on the page.
\end{quote}

\subsection*{What to know first}
\begin{itemize}
\raggedright
  \item \emph{namaste} — you already know \textbf{\hi{न}}, the syllable \emph{te}, and how consonants carry \textquotedblleft{}a.\textquotedblright{}
\end{itemize}

\subsection*{The letters in this word}
\begin{itemize}
\raggedright
  \item \textbf{\hi{म}} = \textquotedblleft{}ma\textquotedblright{} · \textbf{\hi{स}} = \textquotedblleft{}sa.\textquotedblright{}
  \item \textbf{\hi{्}}, the \emph{halant}, \textbf{kills a consonant's vowel}: \textbf{\hi{स्}} = a bare \textquotedblleft{}s,\textquotedblright{} which leans onto the next consonant as a \textbf{conjunct}.
\end{itemize}

So \textbf{\hi{नमस्}} (\emph{namas}, \textquotedblleft{}a bow\textquotedblright{}) reads \textbf{\hi{न}·\hi{म}·\hi{स्}}. Add the last lesson's \emph{te} and every piece of \emph{namaste} is yours. Three facts did it: consonants carry \emph{a}, a \emph{mātrā} changes the vowel, a \emph{halant} removes it.

The second half, \emph{kār}, needs three more shapes; the whole written word waits for the end of this chapter.

\begin{cousinweb}
\emph{Namaskār} = \emph{namaḥ} (\textquotedblleft{}a bow,\textquotedblright{} from root \textbf{\hi{नम्}} \emph{nam}, as in \emph{namaste}) + \emph{kāra} (\textquotedblleft{}the making or doing of\textquotedblright{} — from the Sanskrit root \emph{kṛ}, \textquotedblleft{}to do, make\textquotedblright{}). So \emph{namaskār} is literally \textbf{\textquotedblleft{}the making of a bow,\textquotedblright{} \textquotedblleft{}the act of homage.\textquotedblright{}} Where \emph{namaste} says \textquotedblleft{}bow to \emph{you},\textquotedblright{} \emph{namaskār} names the bow itself — a touch more formal and ceremonial.

That piece \emph{-kāra} (\textquotedblleft{}the making/doing of\textquotedblright{}) is a Sanskrit building block you'll see again — e.g. in \emph{ahaṃkāra} (\textquotedblleft{}the making of 'I'\textquotedblright{} $\to$ ego). Learn it once as \textquotedblleft{}the act of \_\_\_.\textquotedblright{}
\end{cousinweb}

\begin{culture}
\emph{namaskār} is a notch more \textbf{formal} than \emph{namaste} — common in respectful, official, or ceremonial settings, and as a courteous greeting to elders or in writing. Both are \textquotedblleft{}hello\textquotedblright{}; \emph{namaskār} wears a tie.
\end{culture}

\begin{sounds}
New: \hi{म} ma, \hi{स} sa, and the \emph{halant} \hi{्} (\hi{स्} = s). Together \hi{न}·\hi{म}·\hi{स्} spell \emph{namas}. \emph{Read it:} \textbf{\hi{न}·\hi{म}·\hi{स्}} Say \emph{namaskār} in its beats: na · ma · s · kā · r.
\end{sounds}

\subsection*{Your turn}
Say these aloud:

\begin{itemize}
\raggedright
  \item read \textquotedblleft{}\hi{नमस्}\textquotedblright{} — na · ma · s — and what the \emph{halant} under \hi{स} does
  \item the shared root — \emph{nam} (\textquotedblleft{}bow\textquotedblright{}) in both \emph{namaste} and \emph{namaskār}
  \item which is more formal? (\emph{namaskār})
\end{itemize}

\begin{tcolorbox}[breakable,colback=teal!4,colframe=teal!35!black,title={Before you move on}]
What does the piece \emph{-kāra} mean? (\textquotedblleft{}The making/doing of.\textquotedblright{}) So \emph{namaskār} literally names what? (\textquotedblleft{}The making of a bow.\textquotedblright{}) What does a \emph{halant} do to a consonant? (Removes its built-in \emph{a}.)
\end{tcolorbox}
\section[\texorpdfstring{\hi{धन्यवाद} (dhanyavād)}{धन्यवाद (dhanyavād)}]{Dhanyavād — \textquotedblleft{}thank you,\textquotedblright{} the Sanskrit way}

\label{lesson:HI-C01-dhanyavad}



\begin{quote}
\textquotedblleft{}Thank you\textquotedblright{} — and the first hint of a big fact about Hindi: it has \emph{two} vocabularies for many everyday things, one Sanskrit, one Persian. This is the Sanskrit one.
\end{quote}

\subsection*{What to know first}
\begin{itemize}
\raggedright
  \item \emph{namaskār} — you know the \emph{halant}/conjunct idea and how a \emph{mātrā} changes a vowel.
\end{itemize}

\subsection*{The letters in this word}
New today: \textbf{\hi{ध}} = \textquotedblleft{}dha,\textquotedblright{} the breathy \emph{d} that opens the word · \textbf{\hi{व}} = \textquotedblleft{}va\textquotedblright{} · \textbf{\hi{द}} = \textquotedblleft{}da.\textquotedblright{} With a \emph{halant}, the last two spell the Sanskrit root behind the word's second half: \textbf{\hi{वद्}} (\emph{vad}, \textquotedblleft{}to speak\textquotedblright{}).

The rest — a \emph{ya} joined to \textbf{\hi{न}} as a conjunct, and the long \emph{ā} of \emph{vād} — waits, and so does the whole written word, for a later chapter. Today, say it and read \textbf{\hi{ध}} and \textbf{\hi{वद्}}.

\begin{cousinweb}
\emph{Dhanyavād} = \emph{dhanya} (\textquotedblleft{}worthy, blessed, fortunate\textquotedblright{}) + \emph{vāda} (\textquotedblleft{}a saying, an utterance,\textquotedblright{} from the Sanskrit root \textbf{\hi{वद्}} \emph{vad}, \textquotedblleft{}to speak\textquotedblright{}). So \emph{dhanyavād} literally offers \textbf{\textquotedblleft{}an utterance of {[}you are{]} worthy/blessed\textquotedblright{}} — a formal, almost literary way to say thanks. (\emph{dhanya} comes from \emph{dhana}, \textquotedblleft{}wealth\textquotedblright{} — to call someone \emph{dhanya} is to call them enriched, fortunate.)
\end{cousinweb}

\begin{culture}
\emph{dhanyavād} is the \textbf{formal, Sanskritic} \textquotedblleft{}thank you\textquotedblright{} — the one you'd write, or say in respectful or official settings. It can even feel a little \emph{heavy} for small everyday thanks, where Hindi speakers often reach for a lighter, Persian- derived word instead — \emph{shukriyā}, your next lesson. Holding those two side by side is the single best window into how Hindi works.
\end{culture}

\begin{sounds}
New: \hi{ध} dha (a breathy \emph{d}), \hi{व} va, \hi{द} da. Together \hi{व}·\hi{द्} spell \emph{vad}, \textquotedblleft{}to speak.\textquotedblright{} \emph{Read it:} \textbf{\hi{व}·\hi{द्}} Say the whole word in its beats: dha · n · ya · vā · d.
\end{sounds}

\subsection*{Your turn}
Say these aloud:

\begin{itemize}
\raggedright
  \item its beats — dha · n · ya · vā · d $\to$ \textquotedblleft{}dhanyavād\textquotedblright{}
  \item the two pieces — \emph{dhanya} (\textquotedblleft{}worthy\textquotedblright{}) + \emph{vāda} (\textquotedblleft{}saying\textquotedblright{})
  \item read \textquotedblleft{}\hi{ध}\textquotedblright{} and \textquotedblleft{}\hi{वद्}\textquotedblright{} — \emph{dha}, \emph{vad} — and which one carries the breath
\end{itemize}

\begin{tcolorbox}[breakable,colback=teal!4,colframe=teal!35!black,title={Before you move on}]
\emph{Read it:} \textbf{\hi{वद्}} (\emph{vad}, \textquotedblleft{}to speak.\textquotedblright{}) What are the two parts of \emph{dhanyavād}? (\emph{dhanya} \textquotedblleft{}worthy\textquotedblright{} + \emph{vāda} \textquotedblleft{}saying.\textquotedblright{}) Is it the formal or casual \textquotedblleft{}thanks\textquotedblright{}? (Formal, Sanskritic — the casual one is Persian-derived, next.)
\end{tcolorbox}
\section[\texorpdfstring{\hi{शुक्रिया} (shukriyā)}{शुक्रिया (shukriyā)}]{Shukriyā — \textquotedblleft{}thanks,\textquotedblright{} Hindi's other heritage}

\label{lesson:HI-C01-shukriya}



\begin{quote}
The everyday \textquotedblleft{}thanks\textquotedblright{} — and it is \emph{not} a Sanskrit word at all. It came into Hindi from Persian and, before that, Arabic. This lesson is where you meet Hindi's second bloodstream.
\end{quote}

\subsection*{What to know first}
\begin{itemize}
\raggedright
  \item \emph{dhanyavād} — the formal, Sanskrit \textquotedblleft{}thanks\textquotedblright{}; this is its casual counterpart.
\end{itemize}

\subsection*{The letters in this word}
New: \textbf{\hi{क}} = \textquotedblleft{}ka,\textquotedblright{} \textbf{\hi{र}} = \textquotedblleft{}ra,\textquotedblright{} and the vowel sign \textbf{\hi{ि}} for \textbf{\textquotedblleft{}i\textquotedblright{}} (note: \textbf{\hi{ि}} is written \emph{before} the consonant but read \emph{after} it).

\begin{itemize}
\raggedright
  \item \textbf{\hi{क्} + \hi{र} $\to$ \hi{क्र}} = \textquotedblleft{}kra\textquotedblright{} (halant on \hi{क}, joined to \hi{र} — another conjunct).
  \item \textbf{\hi{रि}} = \textquotedblleft{}ri\textquotedblright{} (with the \textbf{\hi{ि}} sign). So the middle of the word, \textbf{\hi{क्रि}}, reads \emph{kri}.
\end{itemize}

The opening \emph{shu} and the closing \emph{yā} need letters you have not met yet, so the whole written word waits for a later chapter. Today, say it and read \textbf{\hi{क्रि}}.

\begin{cousinweb}
\emph{Shukriyā} comes, through Persian (\emph{shukriya}), from \textbf{Arabic} \emph{shukr} ( \textquotedblleft{}gratitude, thanks\textquotedblright{}) — root \textbf{sh-k-r}. It is \emph{the same word} as the Arabic \textbf{shukran} (\textquotedblleft{}thank you\textquotedblright{}): Hindi and Arabic, two utterly different languages and scripts, saying thanks with one Semitic root. No Sanskrit here at all.

\begin{quote}
Why does Hindi have an Arabic word for \textquotedblleft{}thanks\textquotedblright{}? Centuries of Persian as the language of courts and administration in South Asia poured thousands of Persian (and, through Persian, Arabic) words into everyday Hindi/Urdu. So Hindi carries \textbf{two vocabularies}: a Sanskrit layer (\emph{dhanyavād}, \emph{namaste}) and a Perso-Arabic layer (\emph{shukriyā}, and \emph{alvidā} next). Often the Sanskrit word feels formal/literary and the Persian one feels everyday — though both are fully Hindi.
\end{quote}
\end{cousinweb}

\begin{culture}
\emph{shukriyā} is the \textbf{casual, everyday} \textquotedblleft{}thanks\textquotedblright{} — what you'd say to a friend, a shopkeeper, a waiter. \emph{dhanyavād} is its formal twin. Choosing between them is choosing a register — and, quietly, a heritage.
\end{culture}

\begin{sounds}
New: \hi{क} ka, \hi{र} ra, and the sign \hi{ि} for “i” (written \emph{before} the consonant, read \emph{after}). With a conjunct \hi{क्} + \hi{र} $\to$ \hi{क्र} (“kra”), the middle of the word reads \hi{क्रि} $\to$ \emph{kri}.
\end{sounds}

\subsection*{Your turn}
Say these aloud:

\begin{itemize}
\raggedright
  \item its beats — shu · k · ri · yā — then read \textquotedblleft{}\hi{क्रि},\textquotedblright{} \emph{kri}
  \item \textquotedblleft{}shukriyā\textquotedblright{} — and note it shares a root with Arabic \emph{shukran}
  \item formal vs casual \textquotedblleft{}thanks\textquotedblright{} — \emph{dhanyavād} / \emph{shukriyā}
\end{itemize}

\begin{tcolorbox}[breakable,colback=teal!4,colframe=teal!35!black,title={Before you move on}]
Where does \emph{shukriyā} come from — Sanskrit or Persian/Arabic? (Persian, from Arabic \emph{shukr} — same root as Arabic \emph{shukran}.) What are Hindi's two vocabulary layers? (Sanskrit and Perso-Arabic.)
\end{tcolorbox}
\section[\texorpdfstring{\hi{अलविदा} (alvidā)}{अलविदा (alvidā)}]{\hi{अलविदा} (alvidā) — \textquotedblleft{}farewell,\textquotedblright{} and the Arabic \emph{al-} again}

\label{lesson:HI-C01-alvida}



\begin{quote}
A goodbye — and another Perso-Arabic word, carrying a piece of Arabic grammar you may have met from the other side.
\end{quote}

\subsection*{What to know first}
\begin{itemize}
\raggedright
  \item \emph{shukriyā} — Hindi's Persian/Arabic layer; here's another word from it.
\end{itemize}

\subsection*{The letters in this word}
New: the \textbf{independent vowel \hi{अ}} = \textquotedblleft{}a\textquotedblright{} (used when a word \emph{starts} with a vowel — consonants have their built-in \emph{a}, but a standalone one needs its own letter), \textbf{\hi{ल}} = \textquotedblleft{}la,\textquotedblright{} and \textbf{\hi{ा}}, the \emph{mātrā} for long \textbf{\textquotedblleft{}ā\textquotedblright{}}. Reusing \textbf{\hi{व} \hi{द}} and the \textbf{\hi{ि}} sign:

\begin{itemize}
\raggedright
  \item \textbf{\hi{अल}} = \textquotedblleft{}al\textquotedblright{} · \textbf{\hi{वि}} = \textquotedblleft{}vi\textquotedblright{} (with the \textbf{\hi{ि}} sign) · \textbf{\hi{दा}} = \textquotedblleft{}dā\textquotedblright{} (\textbf{\hi{द} + \hi{ा}}).
\end{itemize}

Left to right: \textbf{\hi{अ}·\hi{ल}·\hi{वि}·\hi{दा}} = \emph{a-l-vi-dā} $\to$

\begin{quote}
\textbf{\hi{अलविदा}} = \textbf{alvidā}
\end{quote}

\begin{cousinweb}
\textbf{\hi{अलविदा}} comes through Persian/Urdu from \textbf{Arabic} \emph{al-widāʿ} (\textquotedblleft{}the farewell\textquotedblright{}), from the Arabic root \textbf{w-d-ʿ}, \textquotedblleft{}to leave, to bid farewell.\textquotedblright{} And there at the front is \textbf{\hi{अल}} — the Arabic article \textbf{al-}, \textquotedblleft{}the,\textquotedblright{} the very same \textquotedblleft{}the\textquotedblright{} that rode into English on \emph{algebra} and \emph{alcohol}. So \emph{alvidā} is literally \textbf{\textquotedblleft{}the farewell.\textquotedblright{}}
\end{cousinweb}

\begin{culture}
\emph{alvidā} carries a note of \textbf{finality} — it's the goodbye for a long or final parting, more weighted than an everyday \textquotedblleft{}see you.\textquotedblright{} For casual \textquotedblleft{}bye,\textquotedblright{} Hindi speakers more often use \emph{namaste} again, or the English \textquotedblleft{}bye,\textquotedblright{} or \emph{phir milenge} (\textquotedblleft{}we'll meet again\textquotedblright{}). But \emph{alvidā} is the word when the parting matters — and, like \emph{shukriyā}, it quietly carries centuries of Persian in it.
\end{culture}

\begin{sounds}
New: the independent vowel \hi{अ} (“a,” for a word that \emph{starts} with a vowel), \hi{ल} la, and the sign \hi{ा} for long “ā.” \hi{अ}·\hi{ल}·\hi{वि}·\hi{दा} $\to$ \emph{alvidā}.
\end{sounds}

\subsection*{Your turn}
Say these aloud:

\begin{itemize}
\raggedright
  \item read \textquotedblleft{}\hi{अलविदा}\textquotedblright{} — a · l · vi · dā
  \item the \emph{al-} at the front — the Arabic \textquotedblleft{}the,\textquotedblright{} as in \textquotedblleft{}the farewell\textquotedblright{}
  \item when would you use \emph{alvidā} vs a casual bye? (a weighty/final parting)
\end{itemize}

\begin{tcolorbox}[breakable,colback=teal!4,colframe=teal!35!black,title={Before you move on}]
\emph{alvidā} literally means what, and what's the little \emph{al-} doing there? (\textquotedblleft{}The farewell\textquotedblright{} — \emph{al-} is the Arabic \textquotedblleft{}the.\textquotedblright{}) Sanskrit or Perso-Arabic? (Perso-Arabic.)
\end{tcolorbox}
\section[\texorpdfstring{\hi{शिरोरेखा} (shirorekhā)}{शिरोरेखा (shirorekhā)}]{The line on top — Devanagari's head-line}

\label{lesson:HI-W01-shirorekha-na-ma}



\begin{quote}
Look at any line of Hindi:

\begin{quote}
\hi{नमस्ते}
\end{quote}

Before any individual letter, notice the thing you cannot miss: \textbf{a horizontal line running across the top of the whole word.} Latin letters sit \emph{on} a line. Devanagari letters \textbf{hang from} one.
\end{quote}

\subsection*{You'll want to know — The shirorekhā}
That line is the \textbf{\hi{शिरोरेखा}} (\emph{shirorekhā}) — literally the \textbf{\textquotedblleft{}head-line\textquotedblright{}}: \emph{śiras} \textquotedblleft{}head\textquotedblright{} + \emph{rekhā} \textquotedblleft{}line.\textquotedblright{}

\emph{Śiras} is worth a moment. It descends from Proto-Indo-European *\emph{\'{k}erh\textsubscript{2}-}, \textquotedblleft{}head, horn\textquotedblright{} — the root behind Latin \emph{cornu} (\textquotedblleft{}horn\textquotedblright{}), Latin \emph{cerebrum} (\textquotedblleft{}brain\textquotedblright{}), Greek \emph{kara} (\textquotedblleft{}head\textquotedblright{}), and English \textbf{horn}. The line across the top of Hindi writing is, by ancestry, its \textbf{horn}.

Two habits about it, and they are the opposite of what you'd guess:

\begin{enumerate}
\raggedright
  \item \textbf{You draw it last.} Write the letter's body first, then cap it.
  \item \textbf{Across a word it is one line.} Rather than capping each letter separately, write the letters and then draw a single bar over the whole word. That is why Hindi words look like they're strung on a wire: they are.
\end{enumerate}

A caveat you should have from the first lesson: this is the \textbf{common convention, not a rule}. Plenty of writers cap each letter as they go, and some draw the bar first. Stroke orders in this track are how it is usually taught, not a standard you can be wrong about.

And one note about the word \emph{shirorekhā} itself, written \textbf{\hi{शिरोरेखा}}: it contains a letter — \textbf{\hi{ख}} — that this book does not give a stroke order for. Read the word; don't try to draw it. This track only asks you to hand-write letters whose stroke data actually exists, and it will say so every time that gap shows up.

\subsection*{Your turn}
Keep \textbf{\hi{नमस्ते}} visible. Finger-trace only its horizontal head-line once, from left edge to right edge. Then trace that same visible line once with a pencil. Do not copy the letter bodies yet. Say \textbf{\hi{शिरोरेखा} — head-line} after your hand stops, and compare your traced line with the model immediately.

\begin{tcolorbox}[breakable,colback=teal!4,colframe=teal!35!black,title={Before you move on}]
What is the line across the top called, and what does the word literally mean? (\textbf{Shirorekhā} — \textquotedblleft{}\textbf{head-line}\textquotedblright{}.) When is it usually drawn? (\textbf{Last} — and as \textbf{one} line across the whole word. Usually: it is the common convention, not a rule.) What ancient root sits inside \emph{śiras}, \textquotedblleft{}head\textquotedblright{}? (PIE \emph{\textbf{\'{k}erh\textsubscript{2}-}}, also behind \textbf{horn}.) Next: hang your first two letters from it.
\end{tcolorbox}
\section[\texorpdfstring{\hi{न}, \hi{म} (na, ma)}{न, म (na, ma)}]{\hi{न} and \hi{म} — two bodies under one head-line}

\label{lesson:HI-W01-na-ma}



\begin{quote}
You know the bar is usually drawn last. Now give it two letter bodies to cap.
\end{quote}

\subsection*{You'll want to know — \hi{न} — \textquotedblleft{}na\textquotedblright{}}
Three pieces:

\begin{enumerate}
\raggedright
  \item \textbf{a left bowl} — a rounded scoop on the left.
  \item \textbf{a right spine} — a vertical downstroke on the right.
  \item \textbf{the top bar} — capping both. \emph{Last.}
\end{enumerate}

\subsection*{You'll want to know — \hi{म} — \textquotedblleft{}ma\textquotedblright{}}
Four movements — the same spine, but the body starts at the top:

\begin{enumerate}
\raggedright
  \item \textbf{down and round} — descend, loop, sweep right.
  \item \textbf{up the right spine}, without lifting.
  \item \textbf{back down it.}
  \item \textbf{the top bar} — last, again.
\end{enumerate}

\subsection*{You'll want to know — The shared frame}
Both letters have a \textbf{vertical spine on the right}, with the distinguishing shape hanging to its left. This is Devanagari's commonest skeleton: \textbf{spine on the right, character on the left, bar across the top}. Learn the frame once and many new letters are only a new left-hand part.

\textquotedblleft{}Commonest,\textquotedblright{} not \textquotedblleft{}only\textquotedblright{} — a minority of letters (\textbf{\hi{द}} and \textbf{\hi{र}} among them) have \textbf{no} right-hand spine at all. You will meet them as they come.

\subsection*{Your turn}
\hlblockfigure{\detokenize{figures/HI-W01-na-ma-filmstrip.pdf}}{How these letters are written, one after another, stroke by stroke: \hi{न}, \hi{म}}

Keep \textbf{\hi{न} \hi{म}} visible. Copy \textbf{\hi{न}} once, looking back after any piece you need. Then copy \textbf{\hi{म}} once the same way. Put them side by side and add one head-line across both. Compare each copy with the visible model immediately. Memory is not the task yet; one careful, supported copy is enough.

\begin{tcolorbox}[breakable,colback=teal!4,colframe=teal!35!black,title={Before you move on}]
Draw \textbf{\hi{न}} — its three pieces? (Bowl, spine, bar.) Draw \textbf{\hi{म}} — what makes it different? (It \textbf{starts at the top} and runs into the spine.) Where does the spine sit in both? (\textbf{On the right} — though a minority of letters have none.) Next: the vowel already hiding inside each consonant.
\end{tcolorbox}
\section[\texorpdfstring{\hi{अ} (a)}{अ (a)}]{The vowel that is already there — the abugida}

\label{lesson:HI-W02-abugida-ka-ta}



\begin{quote}
Here is the thing that trips up every English speaker, and it is better to meet it now than to be confused for a month.

\textbf{\hi{क} is not \textquotedblleft{}k\textquotedblright{}. \hi{क} is \textquotedblleft{}ka\textquotedblright{}.}
\end{quote}

\subsection*{You'll want to know — Consonants come with a vowel attached}
In the Latin alphabet, \emph{k} is a bare consonant and you must add a vowel to say anything. In Devanagari, every consonant \textbf{already contains} the vowel \textbf{\textquotedblleft{}a\textquotedblright{}}:

\noindent
\begin{tabularx}{\linewidth}{@{}>{\raggedright\arraybackslash}X>{\raggedright\arraybackslash}X>{\raggedright\arraybackslash}X@{}}
\toprule
\textbf{letter} & \textbf{says} & \textbf{\emph{not}} \\
\midrule
\hi{क} & \textbf{ka} & \textasciitilde{}\textasciitilde{}k\textasciitilde{}\textasciitilde{} \\
\hi{न} & \textbf{na} & \textasciitilde{}\textasciitilde{}n\textasciitilde{}\textasciitilde{} \\
\hi{म} & \textbf{ma} & \textasciitilde{}\textasciitilde{}m\textasciitilde{}\textasciitilde{} \\
\hi{त} & \textbf{ta} & \textasciitilde{}\textasciitilde{}t\textasciitilde{}\textasciitilde{} \\
\bottomrule
\end{tabularx}

That built-in \textquotedblleft{}a\textquotedblright{} is the \textbf{inherent vowel}. You get it free, whether you want it or not — and the next two lessons are entirely about what to do when you \emph{don't} want it.

So the two letters you learned last lesson were never \textquotedblleft{}n\textquotedblright{} and \textquotedblleft{}m.\textquotedblright{} They were \textbf{na} and \textbf{ma} — and if you write them together you have already written a real syllable pair: \textbf{\hi{नम}}, \emph{nama}, the first half of \emph{namaste}.

\subsection*{You'll want to know — What this kind of script is called}
A script that works this way — consonants carrying a default vowel — is an \textbf{abugida}.

The word comes from \textbf{Ge'ez}, the old liturgical language of Ethiopia, whose script works the same way — and it is built more cleverly than \textquotedblleft{}alphabet\textquotedblright{} is.

\emph{ʾä — bu — gi — da} names two things at once:

\begin{itemize}
\raggedright
  \item the \textbf{consonants} ʾ, b, g, d — the first four of the old \textbf{Semitic} letter order (the same order that gives Greek \emph{alpha, beta, gamma, delta});
  \item the \textbf{vowels} ä, u, i, a — the first four of Ge'ez's \textbf{vowel series}.
\end{itemize}

So the word doesn't just recite the start of a list, the way \emph{alphabet} recites \emph{alpha-beta}. It \textbf{demonstrates the system it names}: consonant plus vowel, consonant plus vowel, four times over. A term that performs its own definition.

(Ge'ez's own traditional recitation order is different — \emph{hä-lä-ḥä-mä…} — so \emph{abugida} is a scholar's coinage drawing on the inherited Semitic order, not the Ethiopian schoolroom one.)

\subsection*{Your turn}
\begin{itemize}
\raggedright
  \item \emph{Say it:} \textquotedblleft{}\hi{क} is \textbf{ka}, not k\textquotedblright{}
  \item \emph{Write it:} \hi{नम} — two letters, \textbf{one bar across both} — \textquotedblleft{}nama\textquotedblright{}
  \item \emph{Say it:} \textquotedblleft{}abugida — a consonant plus its built-in vowel\textquotedblright{}
\end{itemize}

\begin{tcolorbox}[breakable,colback=teal!4,colframe=teal!35!black,title={Before you move on}]
What does \hi{क} say on its own? (\textbf{\textquotedblleft{}ka\textquotedblright{}} — the vowel is built in.) What is that vowel called? (The \textbf{inherent vowel}.) What kind of script is this, and where does the word come from? (An \textbf{abugida} — \emph{ʾä-bu-gi-da}, from Ge'ez, naming four \textbf{consonants} of the old Semitic order \textbf{and} four \textbf{vowel series} at once, so the word performs its own definition.) Write \textbf{\hi{नम}} — what does it say? (\emph{Nama}.) Next: two more letters, and the sound-map that orders them.
\end{tcolorbox}
\section[\texorpdfstring{\hi{क}, \hi{त} (ka, ta)}{क, त (ka, ta)}]{\hi{क} and \hi{त} — two letters on a map of the mouth}

\label{lesson:HI-W02-ka-ta-mouth-order}



\begin{quote}
The built-in vowel means these are \textbf{ka} and \textbf{ta}, not bare \emph{k} and \emph{t}. Their shapes share the frame you already know.
\end{quote}

\subsection*{You'll want to know — \hi{क} — \textquotedblleft{}ka\textquotedblright{}}
\begin{enumerate}
\raggedright
  \item \textbf{a left loop}
  \item \textbf{a vertical spine} on the right
  \item \textbf{the top bar} — last.
\end{enumerate}

\subsection*{You'll want to know — \hi{त} — \textquotedblleft{}ta\textquotedblright{}}
\begin{enumerate}
\raggedright
  \item \textbf{a small left curl}
  \item \textbf{a tall spine}
  \item \textbf{the top bar} — last.
\end{enumerate}

Same frame as before: spine right, shape left, bar on top. Only the left-hand part changes.

\subsection*{You'll want to know — Why the letters are in this order}
The heart of the Devanagari consonant order is not a historical accident like A-B-C. Its five stop families (\emph{vargas}) are sorted by \textbf{where in your mouth the closure happens}, starting at the soft palate and walking forward to the lips:

\textbf{\hi{क}} (soft palate) $\to$ \textbf{\hi{च}} (hard palate) $\to$ \textbf{\hi{ट}} (roof, tongue curled back) $\to$ \textbf{\hi{त}} (teeth) $\to$ \textbf{\hi{प}} (lips)

Say those five and feel the closure travel forward. Indian phoneticians organised sounds this way well over two thousand years ago — Pāṇini's work is roughly fourth century BCE — producing a phonetics chart memorised as an alphabet long before Europe attempted the same analysis.

Two honest footnotes: this walk describes the \textbf{five stop families}, not the whole list; and this book does not give a stroke order for \textbf{\hi{ट}}. Read it here, and wait for a source you trust before you write it.

\subsection*{Your turn}
\begin{itemize}
\raggedright
  \item \emph{Write it:} \hi{क} — \textquotedblleft{}loop, spine, bar last\textquotedblright{}
  \item \emph{Write it:} \hi{त} — \textquotedblleft{}curl, spine, bar last\textquotedblright{}
  \item \emph{Say it:} \textquotedblleft{}ka … ca … ṭa … ta … pa\textquotedblright{}, back to front
\end{itemize}

\begin{tcolorbox}[breakable,colback=teal!4,colframe=teal!35!black,title={Before you move on}]
What do \textbf{\hi{क}} and \textbf{\hi{त}} say alone? (\textbf{Ka}, \textbf{ta} — each has an inherent vowel.) What orders the five stop families? (\textbf{Place of articulation}, from the soft palate forward to the lips.) Next: replace that built-in vowel.
\end{tcolorbox}
\section[Practice]{Practice — the Hindi greetings and the first four letters}

\label{lesson:HI-C01-practice}



\begin{quote}
No new words or letters. \emph{Read it:} each aloud, then recall its meaning and which of Hindi's two heritages it comes from
\end{quote}

\begin{grammarlens}[title={What you've built}]
\begin{itemize}
\raggedright
  \item \textbf{\hi{नमस्ते}} \emph{namaste} — \textquotedblleft{}I bow to you\textquotedblright{} (Sanskrit \emph{nam}, \textquotedblleft{}bow\textquotedblright{}); hello \& bye.
  \item \textbf{\hi{नमस्कार}} \emph{namaskār} — \textquotedblleft{}the making of a bow\textquotedblright{} (Sanskrit); more formal.
  \item \emph{dhanyavād} — \textquotedblleft{}thanks\textquotedblright{} (Sanskrit \emph{dhanya} \textquotedblleft{}worthy\textquotedblright{} + \emph{vāda} \textquotedblleft{}saying\textquotedblright{}); formal.
  \item \emph{shukriyā} — \textquotedblleft{}thanks\textquotedblright{} (Persian, $\leftarrow$ Arabic \emph{shukr}); everyday.
  \item \emph{alvidā} — \textquotedblleft{}the farewell\textquotedblright{} (Persian, $\leftarrow$ Arabic \emph{al-widāʿ}); weighty.
\end{itemize}

Two big ideas from one small chapter: \textbf{how Devanagari reads} (consonants carry \emph{a}; a \emph{mātrā} changes the vowel; a \emph{halant} removes it and makes a conjunct), and \textbf{Hindi's double heritage} — a Sanskrit layer (\emph{namaste, dhanyavād}) and a Perso-Arabic layer (\emph{shukriyā, alvidā}), often the formal vs. the everyday word for the same thing.
\end{grammarlens}

\subsection*{Your turn — the first letters, by hand}
\emph{Write it:} \hi{न}, \hi{म}, \hi{क}, \hi{त} — each body first, the bar last — then compare with the shapes in this chapter and repair one

\begin{itemize}
\raggedright
  \item The line on top: its name, and when is it drawn? (\textbf{Shirorekhā}, \textquotedblleft{}head-line\textquotedblright{}; \textbf{last}, one line across the word.)
  \item What frame do the four share? (\textbf{Spine right}, shape left, \textbf{bar on top} — \textbf{\hi{म}} starts at the top, \textbf{\hi{क}} with a loop, \textbf{\hi{त}} with a curl.)
  \item What does \textbf{\hi{क}} say alone? (\textbf{Ka}, not \emph{k}: the \textbf{inherent vowel} is built in — the mark of an \textbf{abugida}, a name from Ge'ez.)
  \item What orders \textbf{\hi{क}} … \textbf{\hi{त}}? (\textbf{Where the mouth closes}, from the soft palate forward to the lips.)
\end{itemize}

\subsection*{Your turn}
Say these aloud:

\begin{itemize}
\raggedright
  \item all five words from memory — the two hellos, the two thanks, and the farewell
  \item greet someone (\emph{namaste}), thank them formally (\emph{dhanyavād}) then casually (\emph{shukriyā}), and take a weighty leave (\emph{alvidā})
  \item for each of the five, Sanskrit or Perso-Arabic?
\end{itemize}

\begin{tcolorbox}[breakable,colback=teal!4,colframe=teal!35!black,title={Before you move on}]
Which \textquotedblleft{}thanks\textquotedblright{} is Sanskrit and which is Persian? (\emph{dhanyavād} / \emph{shukriyā}.) What three facts let you read Devanagari? (Consonants carry \emph{a}; a \emph{mātrā} changes the vowel; a \emph{halant} removes it.) Next chapter: \emph{āp} and \emph{tum} — Hindi's formal and informal \textquotedblleft{}you\textquotedblright{} — and introducing yourself.
\end{tcolorbox}
